\documentclass[11pt,a4paper]{article}
\usepackage[T1]{fontenc}
\usepackage{lmodern}
\usepackage[margin=1in]{geometry}
\usepackage{amsmath,amssymb,amsthm}
\usepackage[hidelinks]{hyperref}
\newtheorem{theorem}{Theorem}
\newtheorem{lemma}[theorem]{Lemma}
\newtheorem{proposition}[theorem]{Proposition}
\newcommand{\R}{\mathbb{R}}
\newcommand{\Reach}{\mathcal{R}}
\newcommand{\dd}{\,\mathrm{d}}
\title{Small-time controllability without polynomial reachable neighborhoods}
\author{Kaiwen Jin\\\small\href{mailto:yui.ui.siki@gmail.com}{yui.ui.siki@gmail.com}}
\date{October 12, 2026}

\begin{document}
\maketitle

\begin{abstract}
We construct a polynomial control-affine system in five dimensions, with two
controls in $[-1,1]$, that is small-time locally controllable at the origin
but admits no polynomial lower bound for its reachable neighborhoods.
Every point on the negative fourth axis reachable in time at most $t$ has
distance less than $t^{11}e^{-1/(6t)}/4$ from the origin.
Nevertheless, exact positive and negative fourth-coordinate loops, followed
by linear steering in the other coordinates, reach a full neighborhood for
each $t>0$. The construction combines a critical quartic identity with a
logarithmic estimate and concentration near its equality profile.
It gives a negative answer to the polynomial-growth question for analytic
systems with bounded measurable controls.
\end{abstract}

\section{The growth question and a counterexample}

Consider a control-affine system near $0\in\R^n$,
\[
 \dot x=X_0(x)+\sum_{j=1}^m u_jX_j(x),\qquad x(0)=0,
 \qquad u\in[-1,1]^m.
\]
Controls are measurable, and solutions are absolutely continuous and satisfy
the equations almost everywhere. Write $\Reach_{\le t}(0)$ for the endpoints
reached in a duration at most $t$. Small-time local controllability at the
origin means that $0$ is an interior point of $\Reach_{\le t}(0)$ for every
sufficiently small $t>0$. A stronger quantitative property is the existence
of $C,T>0$ and an integer $N\ge1$ such that
\begin{equation}\label{eq:polynomial-growth}
 B(0,Ct^N)\subset\Reach_{\le t}(0),\qquad 0<t<T,
\end{equation}
where balls are Euclidean.

Jafarpour's Conjecture~1.2 in the v7 preprint PDF of
\cite{Jafarpour} asks whether analyticity and small-time local
controllability imply this polynomial growth property. The question also
appears as AIM Problem~417 for analytic systems with $n\ge3$ and $X_0(0)=0$
\cite{AIM417}. The following example satisfies these assumptions and violates
\eqref{eq:polynomial-growth}.

\begin{theorem}\label{thm:counterexample}
On $\R^5$, consider
\begin{equation}\label{eq:system}
\begin{aligned}
 \dot x_1&=u, & \dot x_2&=x_1, & \dot x_3&=x_2,\\
 \dot x_4&=x_1^2x_3^2-\left(\frac19+x_5^2\right)x_2^4,
 &\dot x_5&=v,
\end{aligned}
\end{equation}
with measurable controls $|u|,|v|\le1$. Its vector fields are polynomial and
its drift vanishes at the origin. For every $t>0$ there is $\delta(t)>0$ such
that
\begin{equation}\label{eq:qualitative-ball}
 B(0,\delta(t))\subset\Reach_{\le t}(0).
\end{equation}
On the other hand, every endpoint $-r e_4\in\Reach_{\le t}(0)$ with $r>0$
satisfies
\begin{equation}\label{eq:flat-bound}
 r<\frac{t^{11}}4\exp\left(-\frac1{6t}\right).
\end{equation}
Consequently, for every $C,T>0$ and every integer $N\ge1$, there is a
$t\in(0,T)$ for which \eqref{eq:polynomial-growth} fails.
\end{theorem}

The fourth equation is based on the quartic competition studied by
Beauchard and Marbach \cite[Section~6]{BeauchardMarbach}, with a constant
coefficient replaced by $1/9+x_5^2$. Their Lemma~6.1 proves that $1/9$ is a
lower bound for the optimal constant in the associated quartic inequality;
Remark~6.2 leaves its exact value open. We prove the critical identity and
construct concentrating curves directly below. At the coefficient $1/9$,
the fourth displacement of a closed triple-integrator trajectory is
nonnegative. An arbitrarily small positive value of $x_5^2$ permits a
negative displacement, but only after concentration on an extremely small
scale. This produces controllability without a polynomial reachable ball.

Section~\ref{sec:logarithmic} proves the bound for all admissible controls.
Section~\ref{sec:controllability} constructs a full reachable neighborhood
and completes the proof of Theorem~\ref{thm:counterexample}.

\section{A logarithmic estimate for closed trajectories}
\label{sec:logarithmic}

For any finite time interval, the first, fifth, second and third coordinates
of \eqref{eq:system} are obtained by successive integration. The fourth is
then the integral of a bounded polynomial expression. This gives a unique
global absolutely continuous solution for every admissible control and every
initial condition. Changing the initial fourth coordinate simply translates
that coordinate, since none of the equations depends on $x_4$.
From the zero initial state, for $0\le s\le t$,
\begin{equation}\label{eq:trajectory-bounds}
 |x_1(s)|\le s,\qquad |x_2(s)|\le\frac{s^2}{2},\qquad
 |x_3(s)|\le\frac{s^3}{6},\qquad |x_5(s)|\le s.
\end{equation}
In particular, $f=x_3$ is $C^2$, $f''$ is Lipschitz, and $f'''=u$ almost
everywhere.

The estimate needed for the fourth displacement concerns curves with zero
endpoint jets. All integrals in this note are Lebesgue integrals. We use
integration by parts for absolutely continuous functions
\cite[Appendix~3.A, Theorems~3.60 and~3.62]{Hunter}.

\begin{lemma}\label{lem:logarithmic}
Let $t>0$ and let $f\in C^2([0,t])$ have Lipschitz second derivative,
$\|f'''\|_{L^\infty}\le1$, and
\[
 f(0)=f'(0)=f''(0)=f(t)=f'(t)=f''(t)=0.
\]
Set
\[
 A=\int_0^t f'^4\dd s,\qquad
 q=ff''+\frac13f'^2,\qquad E=\int_0^t q^2\dd s.
\]
Then
\begin{equation}\label{eq:critical-identity}
 E=\int_0^t (ff'')^2\dd s-\frac19 A.
\end{equation}
If $A>0$, then
\begin{equation}\label{eq:logarithmic-estimate}
 E\ge\frac{A}{36\log^2(t^9/(4A))}.
\end{equation}
\end{lemma}

\begin{proof}
Integration by parts gives
\[
 \int_0^t ff''f'^2\dd s
 =\frac13\int_0^t f(f'^3)'\dd s=-\frac13 A.
\]
Expanding $q^2$ proves \eqref{eq:critical-identity}.
The initial jets and the third-derivative bound give
\[
 |f''|\le B:=t,\qquad |f'|\le D:=\frac{t^2}{2},\qquad
 |f|\le S:=\frac{t^3}{6}.
\]
Assume $A>0$ and put
\[
 \varepsilon=\frac{A}{6tBD^2}=\frac{2A}{3t^6}.
\]
Since $A\le tD^4=t^9/16$, we have $0<\varepsilon\le t^3/24<S$.
Let $c(z)=\max(-\varepsilon,\min(z,\varepsilon))$.
Integration by parts yields
\begin{equation}\label{eq:low-part}
 A_{\mathrm{lo}}:=\int_{\{|f|<\varepsilon\}}f'^4\dd s
 =-3\int_0^t c(f)f'^2f''\dd s
 \le3\varepsilon tBD^2=\frac A2.
\end{equation}
Here and below, the truncated compositions are Lipschitz and hence
absolutely continuous. For a $C^1$ function $f$, its derivative is zero
almost everywhere on each level set $\{f=a\}$: points on that set with
$f'\ne0$ are isolated in the level set and form a countable set.
At a threshold point with $f'=0$, the chain rule for a Lipschitz composition
follows directly from its Lipschitz bound. Thus the usual chain rules hold
almost everywhere for these finitely truncated functions, and no threshold
level contributes to an integral containing $f'^4$.

Define
\[
 w(z)=
 \begin{cases}
  \log(S/\varepsilon),& |z|\le\varepsilon,\\
  \log(S/|z|),&\varepsilon<|z|<S,\\
  0,&|z|\ge S.
 \end{cases}
\]
This is a bounded Lipschitz function. On the range of $f$,
\[
 0\le w(f)\le L:=\log(S/\varepsilon)=\log\frac{t^9}{4A}.
\]
A second integration by parts gives
\begin{align}
 \int_0^t qf'^2w(f)\dd s
 &=\frac13\int_0^t fw(f)(f'^3)'\dd s
     +\frac13\int_0^t f'^4w(f)\dd s\notag\\
 &=-\frac13\int_0^t fw'(f)f'^4\dd s
   =\frac13 A_{\mathrm{hi}},\label{eq:weighted-identity}
\end{align}
where $A_{\mathrm{hi}}=\int_{\{|f|>\varepsilon\}}f'^4\dd s$.
The boundary terms vanish, and $zw'(z)=-1$ when
$\varepsilon<|z|<S$. The level sets $|f|=\varepsilon$ and $|f|=S$
give zero contributions as just explained.
By \eqref{eq:low-part}, \eqref{eq:weighted-identity} and Cauchy--Schwarz,
\[
 \frac A2\le A_{\mathrm{hi}}
 \le3\sqrt E\left(\int_0^t f'^4w(f)^2\dd s\right)^{1/2}
 \le3\sqrt E\sqrt A L.
\]
Division is legitimate because $A>0$ and $L\ge\log4>0$.
This proves \eqref{eq:logarithmic-estimate}.
\end{proof}

Suppose a control of duration $t>0$ reaches a negative fourth-axis point.
Then $f=x_3$ has all six zero endpoint jets in
Lemma~\ref{lem:logarithmic}. Equations \eqref{eq:system} and
\eqref{eq:critical-identity} give
\begin{equation}\label{eq:fourth-displacement}
 x_4(t)=E-\int_0^t x_5^2f'^4\dd s.
\end{equation}
If $A=0$, continuity gives $f'=0$, and this displacement is zero.
If $x_4(t)<0$, therefore, $A>0$ and $E<t^2A$ by
\eqref{eq:trajectory-bounds}. The logarithmic estimate implies
\[
 \log\frac{t^9}{4A}>\frac1{6t},\qquad
 A<\frac{t^9}{4}e^{-1/(6t)},\qquad
 0<-x_4(t)<\frac{t^{11}}4e^{-1/(6t)}.
\]
The last expression is strictly increasing in $t>0$, so the same bound
holds for every negative fourth-axis point reached in any duration at most
$t$. This proves \eqref{eq:flat-bound} for the entire measurable control
cube, without a restriction on the number of switches.

\section{Concentrating loops and a reachable neighborhood}
\label{sec:controllability}

The equation $ff''+f'^2/3=0$ has the profile $f(s)=s^{3/4}$ on
$s>0$. Its critical energy vanishes, while $f'^4$ has a logarithmically
divergent integral at zero. The following endpoint corrections turn that
profile into admissible closed trajectories.

\begin{lemma}\label{lem:concentration}
There are $C^2$ functions $f_\eta$ on $[0,1]$, with Lipschitz second
derivative and zero endpoint jets through order two, such that
\begin{equation}\label{eq:concentration}
 A_\eta:=\int_0^1 f_\eta'^4\dd s
       =A_0+\frac{81}{256}\log\frac{c}{\eta},\qquad
 E_\eta:=\int_0^1\left(f_\eta f_\eta''+\frac13f_\eta'^2\right)^2\dd s
       =E_0,
\end{equation}
for $0<\eta<c=1/4$, where $A_0,E_0$ are finite positive constants.
In particular, $E_\eta/A_\eta\to0$ as $\eta\downarrow0$.
\end{lemma}

\begin{proof}
Set
\begin{align*}
 h(r)&=\frac{221}{32}r^3-\frac{153}{16}r^4+\frac{117}{32}r^5,\\
 k(r)&=1+\frac34r-\frac3{32}r^2-\frac{455}{32}r^3
             +\frac{663}{32}r^4-\frac{261}{32}r^5.
\end{align*}
Their endpoint jets are
\begin{align*}
 (h,h',h'')(0)&=(0,0,0),& (h,h',h'')(1)&=(1,3/4,-3/16),\\
 (k,k',k'')(0)&=(1,3/4,-3/16),& (k,k',k'')(1)&=(0,0,0).
\end{align*}
For $0<\eta<c$, define
\begin{equation}\label{eq:profile}
 f_\eta(s)=
 \begin{cases}
  \eta^{3/4}h(s/\eta),&0\le s\le\eta,\\
  s^{3/4},&\eta\le s\le c,\\
  c^{3/4}k((s-c)/c),&c\le s\le2c,\\
  0,&2c\le s\le1.
 \end{cases}
\end{equation}
The jets match at all three junctions. Each third derivative is bounded on
its piece for fixed $\eta>0$, and the second derivative is continuous at
the junctions. Consequently $f_\eta''$ is Lipschitz and
$M_\eta:=\|f_\eta'''\|_{L^\infty}<\infty$.

On $[\eta,c]$, the critical expression is zero and
$f_\eta'^4=(81/256)s^{-1}$. In the two polynomial pieces the changes of
variables $r=s/\eta$ and $r=(s-c)/c$ cancel the scale factors exactly.
Thus \eqref{eq:concentration} holds with
\begin{align*}
 A_0&=\int_0^1 h'^4\dd r+\int_0^1 k'^4\dd r,\\
 E_0&=\int_0^1(hh''+h'^2/3)^2\dd r
       +\int_0^1(kk''+k'^2/3)^2\dd r.
\end{align*}
These constants are finite. Also $A_0>0$ since $h'$ is not identically
zero. If $a_0=221/32$, then
$hh''+h'^2/3=9a_0^2r^4+O(r^5)$, so $E_0>0$.
Finally, the logarithm in \eqref{eq:concentration} diverges as
$\eta\downarrow0$.
\end{proof}

\begin{proof}[Proof of Theorem~\ref{thm:counterexample}]
The polynomial vector fields, global solutions and negative-axis estimate
were established in Section~\ref{sec:logarithmic}. We construct a full
neighborhood, fixing an arbitrary $t>0$.
Let $\sigma=t/4$ and $a=t/16$. By
Lemma~\ref{lem:concentration}, choose $\eta_-\in(0,c)$ with
$E_0<a^2A_{\eta_-}$, and set
\[
 \beta_-^{\max}=\frac{\sigma^3}{1+M_{\eta_-}}.
\]
First apply $u=0,v=1$ for a duration $a$.
The first three coordinates remain zero and $x_5=a$; there is no fourth
displacement. Keep $v=0$ during the next interval of length $\sigma$, and
in its local time $s$ use
\[
 u(s)=\beta\sigma^{-3}f_{\eta_-}'''(s/\sigma),\qquad
 0\le\beta\le\beta_-^{\max}.
\]
This measurable control has absolute value at most one. The triple
coordinates are
\[
 (x_1,x_2,x_3)(s)
 =\bigl(\beta\sigma^{-2}f_{\eta_-}''(s/\sigma),
        \beta\sigma^{-1}f_{\eta_-}'(s/\sigma),
        \beta f_{\eta_-}(s/\sigma)\bigr),
\]
so they return to zero. The fourth displacement is exactly
\begin{equation}\label{eq:negative-loop}
 \beta^4\sigma^{-3}(E_0-a^2A_{\eta_-}).
\end{equation}
Apply $u=0,v=-1$ for another duration $a$ to return $x_5$ to zero.
This also gives no fourth displacement. The total duration is $3t/8$,
and varying $\beta$ fills an interval $[-d_-,0]$ with $d_->0$.

For a positive loop take $\eta_+=c/2$ and keep $v=0$.
The same triple trajectory, now with
$0\le\beta\le\beta_+^{\max}:=\sigma^3/(1+M_{\eta_+})$, gives fourth
displacement $\beta^4\sigma^{-3}E_0$ in time $\sigma$.
It fills $[0,d_+]$ for some $d_+>0$.
Both intervals are filled exactly, since every nonnegative real number has
a nonnegative fourth root. The other four coordinates end at zero.
These loops add the same fourth displacement from any initial $x_4$.

It remains to steer the first three and fifth coordinates while compensating
for their fourth displacement. Put $\ell=t/4$.
For the triple integrator from zero, its endpoint under a control $u$ is
\begin{equation}\label{eq:linear-endpoint}
 \left(\int_0^\ell u\dd s,
       \int_0^\ell(\ell-s)u\dd s,
       \int_0^\ell\frac{(\ell-s)^2}{2}u\dd s\right).
\end{equation}
Let $p_0(s)=1$, $p_1(s)=\ell-s$, $p_2(s)=(\ell-s)^2/2$ and
$G_{ij}=\int_0^\ell p_i p_j\dd s$ for $0\le i,j\le2$.
The Gram matrix $G$ is positive definite because these polynomials are
linearly independent. For $z=(z_1,z_2,z_3,z_5)$, choose
\[
 u_z(s)=\sum_{i=0}^2\bigl(G^{-1}(z_1,z_2,z_3)^\mathsf{T}\bigr)_i p_i(s),
 \qquad v_z(s)=\frac{z_5}{\ell}.
\]
They give the endpoint $z$ in the four specified coordinates.
For fixed $\ell$, their sup norms depend continuously on $z$ and vanish
at zero. Thus the controls are admissible for all sufficiently small $z$.
Their fourth displacement $F(z)$ is a polynomial in $z$, with terms of
degrees four and six, because the first three coordinates and $x_5$ depend
linearly on $z$. In particular, $F$ is continuous and $F(0)=0$.

Set $d=\min(d_-,d_+)>0$. Choose $\delta>0$ so small that for
$|y|<\delta$, the controls associated with
$z=(y_1,y_2,y_3,y_5)$ are admissible and
\[
 |y_4-F(z)|<d.
\]
Such a choice follows from continuity at zero, for example by imposing
$|y_4|<d/2$ and $|F(z)|<d/2$.
Execute the scalar loop with displacement $y_4-F(z)$ first, and then
apply $u_z,v_z$ for a duration $\ell$.
The final point is exactly $y$: the second stage starts with the other
four coordinates zero, and its fourth displacement is independent of the
initial $x_4$. The total duration is at most
$3t/8+t/4=5t/8<t$. This proves \eqref{eq:qualitative-ball}.

Finally, write $b(t)=t^{11}e^{-1/(6t)}/4$. For every integer $N\ge1$,
$b(t)/t^N\to0$ as $t\downarrow0$.
Given $C,T>0$, choose $0<t<T$ with $b(t)<Ct^N/2$.
The point $-(C/2)t^N e_4$ lies in $B(0,Ct^N)$ but is excluded from
$\Reach_{\le t}(0)$ by \eqref{eq:flat-bound}.
This proves the failure of every polynomial inner ball and completes the
counterexample.
\end{proof}

\paragraph{AI-use statement.}
OpenAI Codex was used for literature retrieval, mathematical development,
exact arithmetic checks, and manuscript drafting. Independent mathematical
review is requested.

\begin{thebibliography}{9}
\bibitem{Jafarpour}
S.~Jafarpour,
\emph{On small-time local controllability},
SIAM Journal on Control and Optimization \textbf{58} (2020), 425--446.
\href{https://doi.org/10.1137/16M1068797}{doi:10.1137/16M1068797}.
Preprint PDF cited for Conjecture~1.2 (p.~3):
\href{https://arxiv.org/abs/1604.02432v7}{arXiv:1604.02432v7} (2019).

\bibitem{AIM417}
AIM problem repository,
\emph{Polynomial growth of reachable neighborhoods for analytic control systems},
AIM Problem~417,
\href{https://github.com/MColbrook/AIM/blob/c9929805cece044705abaaebe1cebddac1ba11a6/problems/417-analytic-reachable-set-polynomial-growth.md}{versioned statement}.

\bibitem{BeauchardMarbach}
K.~Beauchard and F.~Marbach,
\emph{Some quartic control results for scalar-input systems},
\href{https://arxiv.org/abs/2608.28582v1}{arXiv:2608.28582v1} (2026).

\bibitem{Hunter}
J.~K.~Hunter,
\emph{Notes on Partial Differential Equations}, Appendix~3.A,
\href{https://www.math.ucdavis.edu/~hunter/pdes/ch3A.pdf}{Absolutely continuous functions},
Theorems~3.60 and~3.62, p.~79.
\end{thebibliography}
\end{document}
